
    The input consists of:
    \begin{itemize}
      \item One line with integers $N$, $K$ ($2 \leq N \leq 100$, $1 \leq K \leq 1000$), the length of the permutation and the number of possible periods.

      \item One line with $K$ integers $p_1, p_2, \dots, p_K$ ($1 \leq p_i \leq 10^9$). These are all the possible periods that binary strings of 
      length $N$ can have, with respect to the hidden permutation. The numbers $p_i$ are all distinct, and are sorted in increasing order.

      \item One line with $K$ integers $m_1, m_2, \dots, m_K$ ($1 \leq m_i \leq 2^{100}$). This means 
      that there are $m_i$ binary strings with period $p_i$.

    \end{itemize}
    
    Note that the numbers $m_i$ can be quite large!

